\documentclass[UTF8,12pt,a4paper,fontset=fandol]{ctexbook}


% 支持从项目根目录或当前笔记目录编译。
\makeatletter
\def\input@path{{../}}
\makeatother
\usepackage{researchnotes}


\title{概率论与数理统计笔记}
\author{}
\date{\today}


\begin{document}


\maketitle
\tableofcontents
\clearpage


\chapter*{绪论：研究对象、课程总纲与学习路线}
\addcontentsline{toc}{chapter}{绪论：研究对象、课程总纲与学习路线}

\section*{从随机现象到统计推断}

\keyterm{概率论}（probability theory）研究随机现象的数学模型及其规律；
\keyterm{数理统计}（mathematical statistics）研究如何利用观测数据认识产生数据的总体，并量化结论的不确定性。
两者通过同一套随机模型相互联系：概率论通常从给定模型出发推导观测结果的分布，数理统计则从观测结果出发估计模型中的未知量、检验假设或预测新的观测。

例如，在独立重复的硬币试验中，若已知每次正面出现的概率，可以计算正面总次数的分布；
若这一概率未知，则需要根据有限次试验的结果估计它，并说明估计误差。
本书以这类从模型到数据、再从数据回到模型的问题为主线，逐步建立概念、计算方法与论证依据。

\section*{全书定位与先修知识}

本书按本科阶段的系统课程组织，主干内容覆盖概率基础、随机变量、极限定理、统计推断以及常用统计模型。
读者应具备集合与函数、一元微积分、级数和初步线性代数知识；多元积分、变量替换与矩阵二次型在使用前集中回顾。
测度论不作为起点，但会明确事件族、可测性和可积性的含义，避免把只在离散情形成立的直觉无条件推广到一般情形。

当前版本为\keyterm{课程总纲与分章写作大纲}：各节说明后续正文需要建立的概念、证明的结论、处理的例题与易错点。
其中少量公式用于确定记号和结论的适用范围；完整定义、证明、例题解答和习题解答将在相应章节继续展开。

\section*{课程总纲}

全书分为四篇。前三篇为主干，第四篇为选读；各章顺序反映知识依赖，并非固定课时安排。

\begin{description}
  \item[第一篇：概率模型与随机变量]
  从事件和概率出发，研究分布、数字特征、条件结构与随机变量序列的极限。
  \begin{enumerate}
    \item 第\ref{chap:ps-events}章：随机试验、事件与概率空间。
    \item 第\ref{chap:ps-conditioning}章：条件概率、独立性与贝叶斯公式。
    \item 第\ref{chap:ps-univariate}章：随机变量、分布函数与常见分布。
    \item 第\ref{chap:ps-multivariate}章：多维随机变量及其分布。
    \item 第\ref{chap:ps-moments}章：数学期望与数字特征。
    \item 第\ref{chap:ps-transformations}章：随机变量函数的分布。
    \item 第\ref{chap:ps-conditional-expectation}章：条件期望与条件方差。
    \item 第\ref{chap:ps-generating-functions}章：生成函数与特征函数。
    \item 第\ref{chap:ps-limits}章：概率不等式、大数定律与中心极限定理。
  \end{enumerate}
  \item[第二篇：数理统计的基本理论]
  由随机抽样建立统计模型，依次研究样本包含的信息、点估计、区间估计与假设检验。
  \begin{enumerate}
    \item 第\ref{chap:ps-sampling}章：总体、样本、统计量与抽样分布。
    \item 第\ref{chap:ps-information}章：充分统计量与统计信息。
    \item 第\ref{chap:ps-point-estimation}章：点估计及其评价。
    \item 第\ref{chap:ps-intervals}章：区间估计与置信区域。
    \item 第\ref{chap:ps-testing}章：假设检验的原理与方法。
    \item 第\ref{chap:ps-nonparametric}章：拟合优度、列联表与非参数方法。
  \end{enumerate}
  \item[第三篇：统计模型与数据分析]
  将估计和检验用于多个变量之间的关系与组间差异，形成建模、计算、诊断和解释的完整过程。
  \begin{enumerate}
    \item 第\ref{chap:ps-regression}章：线性回归与预测。
    \item 第\ref{chap:ps-anova}章：方差分析与试验设计初步。
  \end{enumerate}
  \item[第四篇：进一步学习的专题（选读）]
  \begin{enumerate}
    \item 第\ref{chap:ps-bayesian}章：贝叶斯统计初步。
    \item 第\ref{chap:ps-processes}章：随机过程初步。
    \item 第\ref{chap:ps-computation}章：随机模拟与统计计算。
  \end{enumerate}
\end{description}

附录\ref{app:ps-prerequisites}整理必要的数学工具，附录\ref{app:ps-reference}规划分布、推断方法与学习检查索引。
附录用于回顾与查阅，不能代替正文中对条件和推导过程的说明。

\section*{学习路径与章节依赖}

首次学习可沿第一、二篇顺序推进，再学习回归和方差分析。
第\ref{chap:ps-generating-functions}章的特征函数证明细节、
第\ref{chap:ps-information}章的完备性及信息下界、
第\ref{chap:ps-point-estimation}章的渐近理论可在第一遍阅读时暂缓，随后返回补足理论依据。
第\ref{chap:ps-limits}章会区分需要掌握的定理陈述、初等证明以及选读的证明方法。

第\ref{chap:ps-intervals}章和第\ref{chap:ps-testing}章依赖抽样分布与点估计；
第\ref{chap:ps-regression}章和第\ref{chap:ps-anova}章进一步依赖区间估计和检验。
贝叶斯统计以条件概率和似然为基础；随机过程以联合分布、条件期望与极限定理为基础；
统计计算中的独立模拟可在学完极限定理后阅读，自助法在统计推断之后阅读，马尔可夫链蒙特卡罗则在随机过程之后阅读。

每章按“问题引入---定义与条件---性质或定理---例题---习题与小结”扩写。
习题分为概念辨析、基本计算、证明与反例、综合建模四类；计算实验用于辅助理解，并明确其不能替代理论证明。

\section*{全书记号与论述约定}

\begin{description}
  \item[概率模型]
  $\Omega$ 表示样本空间，$\mathcal F$ 表示事件族，$\mathbb P$ 表示概率；事件通常记为 $A,B$。
  这些对象在第\ref{chap:ps-events}章正式定义。
  \item[随机变量与观测]
  $X,Y$ 及 $X_1,\ldots,X_n$ 表示随机变量；$x,y$ 及 $x_1,\ldots,x_n$ 表示其取值或观测值。
  $n$ 表示样本量，$\mathbf X=(X_1,\ldots,X_n)$ 表示样本向量，$\mathbf x$ 表示其观测值。
  \item[分布与数字特征]
  $F_X$、$p_X$ 和 $f_X$ 分别表示分布函数、概率质量函数和概率密度函数。
  $\mathbb E[X]$、$\operatorname{Var}(X)$ 和 $\operatorname{Cov}(X,Y)$ 分别表示期望、方差和协方差；使用前说明其存在条件。
  \item[统计模型与估计]
  $\theta$ 表示未知参数，$\Theta$ 表示参数空间，$\mathbb P_\theta$ 和 $\mathbb E_\theta$ 表示参数为 $\theta$ 时的概率与期望。
  $T(\mathbf X)$ 表示统计量；$\widehat\theta(\mathbf X)$ 表示估计量，$\widehat\theta(\mathbf x)$ 表示观测后得到的估计值。
  \item[假设与结论]
  独立性、同分布、正态性和同方差性均是需要说明的模型条件。
  精确结论、渐近结论和数值近似分别表述；不把大样本近似写成有限样本恒等式。
\end{description}


\part{概率模型与随机变量}

\chapter{随机试验、事件与概率空间}
\label{chap:ps-events}

本章从“怎样把一次随机试验写成数学对象”出发，建立概率论的基本语言。
学习目标是正确选择样本空间、表达事件，并说明所采用的概率模型为什么适合问题。先修知识为集合运算与初等计数。

\section{随机试验与事件的集合表示}
\label{sec:ps-events-language}
引入\keyterm{随机试验}（random experiment）、\keyterm{样本空间}（sample space）与\keyterm{事件}（event）。
以抛硬币、抽球和等待时间说明有限、可数与不可数样本空间；讲解事件的并、交、补、差以及德摩根律。
区分互斥事件、对立事件、必然事件与不可能事件，并训练把“至少”“至多”“恰好”翻译为集合关系。

\section{古典概型与计数方法}
\label{sec:ps-counting}
在有限样本空间且基本结果等可能的条件下建立\keyterm{古典概型}（classical probability model）。
整理加法原理、乘法原理、排列、组合与有限容斥公式；比较有放回和无放回、有序和无序抽样。
先验证等可能性，再使用有利结果数与总结果数之比，避免对未经论证的结果分类直接等权计数。

\section{概率公理与基本性质}
\label{sec:ps-probability-axioms}
定义\keyterm{$\sigma$ 代数}（sigma-algebra）$\mathcal F$ 与\keyterm{概率空间}（probability space）$(\Omega,\mathcal F,\mathbb P)$，
说明非负性、规范性与可数可加性。
由公理证明补事件公式、单调性、有限加法公式、并集上界以及概率对单调事件列的连续性；强调一般样本空间中事件属于指定的事件族。

\section{几何概率与建模边界}
\label{sec:ps-geometric-probability}
在有限且非零的长度、面积或体积区域上，说明均匀分布假设如何导出\keyterm{几何概率}（geometric probability）。
比较不同随机抽取机制产生的不同模型，区分概率为零与不可能、概率为一与逻辑必然。
为连续型随机变量中“单点概率为零”预留直观基础。

\section{例题、习题与学习检查}
规划生日碰撞、无放回抽样、随机到达的会面问题三个例题，逐项写明试验机制、样本空间与目标事件。
证明题训练从公理推导结论；辨析题要求构造“各类结果不等可能”的反例。
完成本章后，应能说明一个概率答案所依赖的模型，而不只给出计数结果。


\chapter{条件概率、独立性与贝叶斯公式}
\label{chap:ps-conditioning}

本章研究获得新信息后概率如何变化，以及何时一个事件不会改变另一个事件的概率。
先修内容为第\ref{chap:ps-events}章；学习目标是正确选择条件、划分情形并辨认独立性假设。

\section{条件概率与乘法公式}
\label{sec:ps-conditional-probability}
定义\keyterm{条件概率}（conditional probability）$\mathbb P(A\mid B)$，首先限定 $\mathbb P(B)>0$。
推导乘法公式与多步试验的链式分解，说明树状分析的每条分支都带有相应条件。
零概率事件不能直接代入比值定义，其连续条件分布处理留待第\ref{chap:ps-multivariate}章。

\section{全概率公式与贝叶斯公式}
\label{sec:ps-total-bayes}
在有限或可数的事件划分下建立\keyterm{全概率公式}（law of total probability）与\keyterm{贝叶斯公式}（Bayes' formula），
明确互斥、覆盖与相关分母为正的条件。
由分层抽样和检测模型说明先验信息、条件概率与后验信息的区别，避免颠倒条件。

\section{事件的独立性}
\label{sec:ps-event-independence}
定义\keyterm{独立性}（independence）、两两独立和相互独立；对有限事件族要求相应各子族满足乘法关系。
证明独立性在取补运算下的保持性质，构造两两独立但不相互独立的例子。
比较互斥与独立，并说明非零概率的互斥事件不能独立。

\section{条件独立与条件化策略}
\label{sec:ps-conditional-independence}
定义给定正概率事件后的\keyterm{条件独立}（conditional independence），说明它与无条件独立之间一般不能相互推出。
介绍按第一步、来源类别或隐藏状态条件化的解题方法，并用分层比例说明总体关系与分组关系可能不同。

\section{例题、习题与学习检查}
规划双袋抽球、含假阳性与假阴性的检测、简单系统可靠性、两两独立反例等例题；检测例题使用明确给定的假想参数。
习题要求分别写出模型假设和推导步骤，并用直接计数核对小规模情形。
学习检查聚焦“条件是什么”“划分是否完备”“独立性从何而来”三个问题。


\chapter{随机变量、分布函数与常见分布}
\label{chap:ps-univariate}

本章把事件模型转化为数值模型，建立描述单个随机量的统一方法。
先修内容为事件与条件概率；密度部分需要一元积分。学习目标是从试验机制识别分布，并按分布计算事件概率。

\section{随机变量与分布函数}
\label{sec:ps-random-variable-cdf}
定义实值\keyterm{随机变量}（random variable）$X$，说明对每个实数 $x$，集合 $\{X\leq x\}$ 都必须是事件。
定义\keyterm{分布函数}（cumulative distribution function, CDF）$F_X(x)=\mathbb P(X\leq x)$，证明其单调性、右连续性和两端极限性质。
由分布函数的跳跃计算单点概率，并用左极限处理开闭端点。

\section{离散型随机变量与常见计数模型}
\label{sec:ps-discrete-distributions}
定义\keyterm{概率质量函数}（probability mass function, PMF），说明非负性与总和为一。
依次介绍伯努利分布（Bernoulli）、二项分布（binomial）、超几何分布（hypergeometric）与泊松分布（Poisson），列出参数范围、取值范围和试验机制。
比较固定次数独立试验、无放回抽样与稀少事件计数；泊松近似的极限条件在第\ref{chap:ps-limits}章讨论。

\section{等待型离散分布与记号约定}
\label{sec:ps-discrete-waiting}
介绍几何分布（geometric）与负二项分布（negative binomial），统一以“获得指定成功次数所需的总试验次数”为基本版本。
另列“成功前的失败次数”版本及其平移关系，防止取值起点和参数化混用。
由独立伯努利试验推导质量函数，并证明几何分布的无记忆性。

\section{密度与常见连续分布}
\label{sec:ps-continuous-distributions}
定义\keyterm{概率密度函数}（probability density function, PDF），明确本书的连续型计算主要指具有密度的绝对连续分布。
依次介绍均匀分布（uniform）、指数分布（exponential）、正态分布（normal）、伽马分布（gamma）与贝塔分布（beta）。
伽马分布统一采用形状与率参数化，同时给出尺度与率的倒数关系；密度值不直接等于单点概率。

\section{分位数、混合分布与分布的边界}
\label{sec:ps-quantiles-mixtures}
对 $0<u<1$ 定义\keyterm{分位数}（quantile）$q_u=\inf\{x:F_X(x)\geq u\}$，说明非严格单调分布函数仍可使用广义逆。
用“零点有质量、正半轴有密度”的模型介绍\keyterm{混合分布}（mixed distribution）。
指出分布函数连续并不必然意味着存在密度，离散与具有密度两类也不穷尽一切分布。

\section{例题、习题与学习检查}
规划按分段分布函数求概率、按抽样机制选择二项或超几何模型、按等待机制选择几何或指数模型、正态标准化与分位数计算。
习题要求先检查质量函数或密度是否归一化，再计算；反例题区分“同一分布”与“同一个随机变量”。
本章暂不集中计算均值和方差，相关数字特征在第\ref{chap:ps-moments}章统一建立。


\chapter{多维随机变量及其分布}
\label{chap:ps-multivariate}

本章研究多个随机量之间的联合关系，回答“分别知道各自的分布是否足够”的问题。
先修内容为第\ref{chap:ps-univariate}章及二重积分；学习目标是掌握联合、边缘、条件分布与独立性之间的联系。

\section{随机向量与联合分布}
\label{sec:ps-joint-distributions}
定义\keyterm{随机向量}（random vector）与\keyterm{联合分布}（joint distribution）。
以二维情形建立联合分布函数、联合质量函数与联合密度，明确分布的支撑区域。
训练通过事件区域求和或积分，避免只写积分核而遗漏区域约束。

\section{边缘分布与条件分布}
\label{sec:ps-marginal-conditional}
由联合分布推出\keyterm{边缘分布}（marginal distribution）与\keyterm{条件分布}（conditional distribution）。
离散情形要求被条件化取值具有正概率；具有联合密度时，在边缘密度为正处给出条件密度公式。
解释连续条件分布并非把两个零概率事件直接相除，并说明条件分布在零概率取值上的版本可以不唯一。

\section{随机变量的独立性}
\label{sec:ps-variable-independence}
用联合分布对边缘分布的分解定义随机变量的独立性，给出分布函数及质量函数、密度的相应判据。
说明密度分解按几乎处处理解，检查乘积分解时必须包含支撑区域。
证明独立随机变量分别作可测函数变换后仍独立，并区分两两独立与相互独立。

\section{常见多维分布}
\label{sec:ps-multivariate-models}
介绍多项分布（multinomial distribution），由互斥类别的独立重复试验导出其质量函数。
介绍二维正态分布（bivariate normal distribution）的参数与密度结构；一般多元正态模型在矩阵工具齐备后补充。
指出两个正态边缘分布不保证联合正态，也不保证独立；协方差判据在第\ref{chap:ps-moments}章建立。

\section{例题、习题与学习检查}
规划三角形区域上的均匀分布、带约束的联合质量表、给定总次数的类别计数等例题。
习题要求从同一联合模型计算边缘和条件概率，并构造边缘相同而联合关系不同的两个模型。
学习检查包括区域边界、归一化、条件分母及独立性依据。


\chapter{数学期望与数字特征}
\label{chap:ps-moments}

本章研究怎样用少量数值描述分布，以及这些数值遗漏了什么信息。
先修内容为单变量与联合分布；学习目标是正确计算期望、方差和协方差，并在计算前检查存在条件。

\section{期望的定义与存在性}
\label{sec:ps-expectation-definition}
定义\keyterm{数学期望}（mathematical expectation），由离散加权和与密度积分建立统一认识。
区分非负随机变量的扩展期望、一般随机变量的期望存在与有限可积；有限期望采用 $\mathbb E[|X|]<\infty$ 作为条件。
利用重尾例子说明对称性不能替代可积性，不能用发散正负项的抵消定义期望。

\section{期望的计算与基本性质}
\label{sec:ps-expectation-properties}
建立随机变量函数的期望公式、线性性、单调性以及指示变量方法。
说明有限和的期望线性性不要求独立，而乘积期望分解需要独立性及相应可积条件。
对无限求和、积分与期望的交换分别列明非负性或绝对可积性等条件。

\section{方差、矩与位置特征}
\label{sec:ps-variance-moments}
定义\keyterm{方差}（variance）、标准差、原点矩与中心矩，说明方差有限的二阶矩条件。
推导平移、缩放与方差计算公式；比较均值、中位数和众数的含义。
偏度和峰度作为选读内容，分别注明所需矩存在以及峰度是否减去 $3$ 的约定。

\section{协方差、相关与随机向量的数字特征}
\label{sec:ps-covariance-correlation}
在二阶矩有限条件下定义\keyterm{协方差}（covariance）与\keyterm{相关系数}（correlation coefficient），相关系数还要求两个方差均非零。
推导随机变量和的方差，借助柯西--施瓦茨不等式（Cauchy--Schwarz inequality）说明相关系数的范围，并构造不相关但不独立的例子。
引入均值向量与协方差矩阵，证明后者半正定；说明联合正态情形下不相关可推出独立。

\section{例题、习题与学习检查}
规划用指示变量求匹配次数期望、比较有放回和无放回抽样的方差、构造零相关的非线性依赖、检查重尾分布的矩。
统一计算第\ref{chap:ps-univariate}章常见分布的均值和方差，记录参数化及成立条件。
习题要求说明所用性质是否依赖独立性，区分“期望相同”与“分布相同”。


\chapter{随机变量函数的分布}
\label{chap:ps-transformations}

本章研究由已知随机量构造的新随机量，其分布如何求得。
先修内容为联合分布与微积分；学习目标是根据变换性质选择分布函数法、变量替换法、卷积或条件化。

\section{分布函数法与离散变换}
\label{sec:ps-cdf-transform}
从事件 $\{g(X)\leq y\}$ 出发求变换后的分布，明确函数 $g$、原变量 $X$ 与新变量 $Y=g(X)$ 的角色。
离散情形对相同函数值合并概率；连续情形处理单调变换、平方与绝对值等非单调变换。
优先确定新变量的取值范围，并保留可能存在的原子概率。

\section{密度变换与雅可比行列式}
\label{sec:ps-jacobian}
在适当的可微、可逆及雅可比非退化条件下建立一维和多维密度变换公式。
说明使用逆变换的\keyterm{雅可比行列式}（Jacobian determinant）绝对值，并将原支撑区域变换到新坐标中。
多分支情形分区求和；用极坐标和商的分布展示方法的适用范围。

\section{和、差及其他组合的分布}
\label{sec:ps-convolution}
建立独立随机变量和的\keyterm{卷积}（convolution）公式，比较离散求和与连续积分。
推导独立同成功概率二项变量、独立泊松变量、独立正态变量以及独立同率参数伽马变量的可加性。
对于相关变量，回到联合分布或条件分布计算，不能直接套用独立卷积。

\section{最大值、最小值与次序统计量的准备}
\label{sec:ps-extreme-transform}
对独立样本用事件交并求最大值和最小值的分布，比较同分布与不同分布情形。
介绍样本排序的含义，为第\ref{chap:ps-sampling}章的次序统计量作准备。
用极值与均值的不同归一化说明不同统计量可能具有不同的误差尺度。

\section{例题、习题与学习检查}
规划正态变量平方、独立均匀变量之和、两个独立指数变量的最小值、连续变量之比等例题。
每题按“支撑区域---变换关系---分布计算---归一化核对”组织。
习题安排同一问题的两种解法，并检查是否遗漏非单调变换的分支。


\chapter{条件期望与条件方差}
\label{chap:ps-conditional-expectation}

本章将“利用已知信息重新求平均”发展为系统工具。
先修内容为条件分布与数学期望；学习目标是理解条件期望本身可以是随机变量，并用它分解复杂计算。

\section{给定事件和随机变量的条件期望}
\label{sec:ps-conditional-expectation-definition}
由条件分布定义\keyterm{条件期望}（conditional expectation）。
区分固定取值 $y$ 下的数值 $\mathbb E[X\mid Y=y]$ 与随机变量 $\mathbb E[X\mid Y]$；以 $X$ 可积作为基本条件。
连续情形说明定义按条件分布的版本理解，不能把“给定取值”机械替换成正概率事件的比值。

\section{全期望公式与迭代性质}
\label{sec:ps-tower-property}
建立\keyterm{全期望公式}（law of total expectation）、线性性、已知量提出与信息嵌套下的迭代性质。
使用 $\sigma$ 代数表示信息的严格定义作为选读，明确可测性、积分刻画与几乎处处唯一性。
通过嵌套与非嵌套信息的比较说明塔式性质不能随意交换条件。

\section{条件方差与全方差公式}
\label{sec:ps-total-variance}
在 $X$ 二阶矩有限时定义\keyterm{条件方差}（conditional variance），建立\keyterm{全方差公式}（law of total variance）：
\begin{equation}
  \operatorname{Var}(X)
  =\mathbb E[\operatorname{Var}(X\mid Y)]
  +\operatorname{Var}(\mathbb E[X\mid Y]).
  \label{eq:ps-total-variance}
\end{equation}
解释组内波动与组间波动的含义，并通过中心化展开给出后续正文中的证明路线。

\section{条件期望作为最优平方损失预测}
\label{sec:ps-conditional-prediction}
在平方可积条件下，研究用 $Y$ 的可测函数预测 $X$ 的问题，证明条件期望使总体均方误差最小。
从误差正交分解解释“最优”的候选范围与几乎处处唯一性。
区分总体最优预测函数与由有限数据估计的回归函数，为第\ref{chap:ps-regression}章作铺垫。

\section{例题、习题与学习检查}
规划随机次数求和、分层总体均值与方差、给定联合密度求条件均值等例题。
随机和例题明确计数变量与被求和序列的独立性及所需矩条件。
习题重点检查“结果是数还是随机变量”、条件信息是否嵌套，以及平方损失结论的适用范围。


\chapter{生成函数与特征函数}
\label{chap:ps-generating-functions}

本章研究如何把分布编码为函数，以便处理矩、独立和与极限。
先修内容为期望、幂级数和独立性；特征函数部分另需复数与欧拉公式。首次学习可先掌握基本计算，再阅读唯一性与收敛理论。

\section{概率生成函数}
\label{sec:ps-pgf}
对非负整数值随机变量定义\keyterm{概率生成函数}（probability generating function, PGF）。
说明在适当收敛区域内由系数恢复概率、由导数求阶乘矩，以及独立和对应函数乘积。
边界处求导需单独检查所需矩是否有限。

\section{矩母函数}
\label{sec:ps-mgf}
定义\keyterm{矩母函数}（moment-generating function, MGF）$M_X(t)=\mathbb E[e^{tX}]$，其中 $t$ 为实数。
区分在单个点有限与在零点的某个开邻域内有限；在后者条件下讨论求导取矩及分布唯一性。
以具有有限若干矩但矩母函数不存在于零点邻域的分布说明方法边界。

\section{特征函数及其基本性质}
\label{sec:ps-characteristic-function}
令 $i$ 为虚数单位，定义\keyterm{特征函数}（characteristic function）$\varphi_X(t)=\mathbb E[e^{itX}]$。
说明它对所有实数 $t$ 都存在，推导平移、缩放和独立和的公式。
陈述特征函数的唯一性定理；连续性定理作为第\ref{chap:ps-limits}章分布收敛证明的选读工具，明确极限函数在零点连续的条件。

\section{例题、习题与学习检查}
规划伯努利、泊松、正态与伽马分布的变换计算，以及独立和分布的识别。
习题比较“所有矩存在”“矩母函数在零点邻域有限”“矩能唯一确定分布”这些不同命题，不将它们混同。
学习检查要求说明每次求导与期望交换的依据，并选择最适合当前分布的工具。


\chapter{概率不等式、大数定律与中心极限定理}
\label{chap:ps-limits}

本章解释样本平均为什么能够反映总体，以及误差为何常用正态分布近似。
先修内容为期望、方差与随机变量序列；学习目标是区分误差界、收敛类型和极限分布，并逐项核对定理条件。

\section{基本概率不等式}
\label{sec:ps-probability-inequalities}
证明马尔可夫不等式（Markov's inequality）与切比雪夫不等式（Chebyshev's inequality），分别说明非负性、阈值为正和二阶矩条件。
系统整理柯西--施瓦茨不等式（Cauchy--Schwarz inequality）及詹森不等式（Jensen's inequality）的概率形式与可积条件。
切尔诺夫界（Chernoff bound）和霍夫丁不等式（Hoeffding's inequality）作为选读，分别检查指数矩与独立有界条件。

\section{随机变量序列的收敛方式}
\label{sec:ps-convergence-modes}
定义\keyterm{几乎必然收敛}（almost sure convergence）、\keyterm{依概率收敛}（convergence in probability）、
\keyterm{依分布收敛}（convergence in distribution）及\keyterm{均方收敛}（mean-square convergence）。
建立常用蕴含关系并配反例；依分布收敛在极限分布函数的连续点定义。
说明依分布收敛到常数可推出依概率收敛，但依概率收敛本身不保证期望收敛。

\section{弱大数定律与强大数定律}
\label{sec:ps-laws-large-numbers}
对独立同分布随机变量 $X_1,X_2,\ldots$，记 $\overline X_n=n^{-1}\sum_{j=1}^nX_j$。
先在共同方差有限时用切比雪夫不等式证明\keyterm{弱大数定律}（weak law of large numbers）。
再陈述独立同分布且 $\mathbb E[|X_1|]<\infty$ 时的\keyterm{强大数定律}（strong law of large numbers），其完整证明列为选读。
区分依概率与几乎必然两种结论，不把一次样本轨迹的单调改善视为定理内容。

\section{中心极限定理与正态近似}
\label{sec:ps-central-limit}
对独立同分布、均值为 $\mu$、方差为 $\sigma^2\in(0,\infty)$ 的随机变量，陈述\keyterm{中心极限定理}（central limit theorem）：
\begin{equation}
  \frac{\sqrt n(\overline X_n-\mu)}{\sigma}
  \xrightarrow{\mathrm d} N(0,1),
  \label{eq:ps-clt}
\end{equation}
其中 $\xrightarrow{\mathrm d}$ 表示依分布收敛，$N(0,1)$ 表示标准正态分布。
用特征函数给出选读证明，讨论二项分布的正态近似与连续性修正。
强调近似质量与原分布及目标尾部有关，不设一个对所有分布都可靠的固定样本量阈值。

\section{连续映射、斯卢茨基定理与德尔塔方法}
\label{sec:ps-asymptotic-tools}
建立\keyterm{连续映射定理}（continuous mapping theorem）与\keyterm{斯卢茨基定理}（Slutsky's theorem），明确函数连续性与另一个因子依概率趋于常数等条件。
在渐近正态与可微条件下介绍\keyterm{德尔塔方法}（delta method），非退化的一阶近似另要求相关导数非零。
由此解释用一致估计的标准差替代总体标准差的依据，为统计推断作准备。

\section{其他近似与适用范围}
\label{sec:ps-approximation-boundaries}
讨论二项分布在试验次数趋于无穷、成功概率趋于零且其乘积趋于有限正数时的泊松极限。
通过有限总体抽样、依赖序列与无限方差分布，说明不能直接套用独立同分布的有限方差中心极限定理。
比较有限样本概率界与渐近近似：前者给可核查的界，后者描述适当极限下的分布形状。

\section{例题、习题与学习检查}
规划用切比雪夫不等式控制样本均值误差、用中心极限定理近似计数概率、构造依概率收敛但期望不收敛的例子。
证明题训练各收敛方式的关系；应用题要求逐条写明独立性、同分布和矩条件。
学习检查以区分“大数定律描述趋近目标”“中心极限定理描述标准化波动”为主。


\part{数理统计的基本理论}

\chapter{总体、样本、统计量与抽样分布}
\label{chap:ps-sampling}

本章从已知概率模型转向未知总体的学习问题，为所有推断方法建立共同框架。
先修内容为分布、数字特征与极限定理；学习目标是区分总体量、样本统计量及其抽样分布。

\section{总体、统计模型与抽样机制}
\label{sec:ps-statistical-model}
定义\keyterm{总体}（population）、\keyterm{统计模型}（statistical model）$\{\mathbb P_\theta:\theta\in\Theta\}$ 与参数空间。
本书将来自同一分布的独立同分布样本称为\keyterm{简单随机样本}（independent and identically distributed sample），并说明有限总体无放回简单随机抽样是另一种机制，其观测通常不独立。
区分随机抽样、随机分组与便利取样，说明推断范围受数据获得方式约束。

\section{统计量与经验分布}
\label{sec:ps-statistics-empirical}
定义不含未知参数的\keyterm{统计量}（statistic），比较统计量、估计量和估计值。
介绍\keyterm{经验分布函数}（empirical distribution function），以及直方图、箱线图和分位数等描述方式。
说明描述数据的图表与关于总体的推断有不同任务，分组和汇总可能掩盖数据结构。

\section{样本均值、样本方差与样本矩}
\label{sec:ps-sample-moments}
对 $n\geq2$ 的样本定义样本均值 $\overline X_n$ 与样本方差
\begin{equation}
  S^2=\frac{1}{n-1}\sum_{j=1}^n(X_j-\overline X_n)^2.
  \label{eq:ps-sample-variance}
\end{equation}
将估计量的期望等于目标总体量称为\keyterm{无偏性}（unbiasedness）。
在独立同分布且二阶矩有限时，证明样本均值和样本方差分别无偏估计总体均值和总体方差，解释分母 $n-1$ 的来源。
比较样本方差、样本二阶中心矩和样本均值的方差，避免把它们当作同一对象。

\section{正态总体的抽样分布}
\label{sec:ps-normal-sampling}
定义卡方分布（chi-square）、$t$ 分布（Student's $t$）与 $F$ 分布，并说明构造中标准正态变量或卡方变量的独立性要求。
对正态独立同分布样本建立均值与样本方差的分布及独立性；推导单样本 $t$ 统计量和独立正态样本的方差比。
自由度的含义通过独立信息与线性约束解释，相关正交分解证明按所需线性代数程度展开。

\section{次序统计量与样本分位数}
\label{sec:ps-order-statistics}
定义\keyterm{次序统计量}（order statistic）$X_{(1)}\leq\cdots\leq X_{(n)}$。
对独立同分布且具有密度的样本，推导第 $k$ 个次序统计量的密度，并说明 $1\leq k\leq n$。
研究极差、中位数和样本分位数，区分离散并列值与连续无并列值情形。

\section{例题、习题与学习检查}
规划由正态样本构造 $t$ 统计量、求均匀样本最大值分布、比较有偏与无偏样本方差公式。
习题安排“是否为统计量”的辨析、抽样分布推导与抽样偏差讨论。
学习检查要求说明重复抽样中变化的对象，并区分样本标准差和估计量的标准误。


\chapter{充分统计量与统计信息}
\label{chap:ps-information}

本章研究能否压缩样本而保留关于参数的信息，以及如何据此改进估计。
先修内容为统计模型、条件分布和条件期望。充分性为主干，完备性与信息下界可在掌握点估计之后回读。

\section{似然函数与因子分解}
\label{sec:ps-likelihood-sufficiency}
定义\keyterm{似然函数}（likelihood function）：固定观测 $\mathbf x$，将联合质量函数或密度视为参数 $\theta$ 的函数。
定义\keyterm{充分统计量}（sufficient statistic），从给定统计量后的样本条件分布不再依赖参数解释其含义。
在共同支配测度下陈述因子分解定理，强调参数依赖的支撑限制属于似然的一部分。

\section{充分性的例子与最小充分性}
\label{sec:ps-minimal-sufficiency}
研究伯努利样本总和、泊松样本总和以及正态模型在不同已知参数条件下的充分统计量。
通过参数影响支撑端点的均匀模型，说明只关注指数项或代数表达式会遗漏信息。
\keyterm{最小充分性}（minimal sufficiency）及其判据列为选读，明确适用模型与零密度情形。

\section{条件化改进、完备性与无偏估计}
\label{sec:ps-rao-blackwell}
在平方可积无偏估计量与充分统计量的条件下，证明\keyterm{拉奥--布莱克韦尔定理}（Rao--Blackwell theorem）的方差改进结论。
定义\keyterm{完备性}（completeness），再陈述莱曼--谢费定理（Lehmann--Scheffé theorem），说明完备充分统计量的无偏函数如何给出唯一的最小方差无偏估计。
将唯一性理解为模型下的几乎处处唯一，并区分充分性与完备性。

\section{费舍尔信息与信息下界}
\label{sec:ps-fisher-information}
定义得分函数（score function）与\keyterm{费舍尔信息}（Fisher information）。
在可微、积分与微分可交换、得分二阶矩有限等正则条件下，建立信息的常用表示与独立样本的可加性。
陈述克拉默--拉奥下界（Cramér--Rao lower bound）的无偏估计版本；用参数相关支撑的例子说明不能越过正则条件。

\section{例题、习题与学习检查}
规划伯努利概率估计的条件化改进、正态均值的信息计算、均匀模型中常规信息下界的适用性分析。
习题要求区分“统计量包含信息”与“估计值接近真值”，并检查改进结论比较的是哪一类估计量。
完整应用将在第\ref{chap:ps-point-estimation}章与估计方法结合。


\chapter{点估计及其评价}
\label{chap:ps-point-estimation}

本章将未知参数估计为一个由样本计算的数值，并建立比较估计方法的标准。
先修内容为抽样分布与似然；涉及最优无偏估计时使用第\ref{chap:ps-information}章，涉及大样本性质时使用第\ref{chap:ps-limits}章。

\section{估计目标、损失与风险}
\label{sec:ps-estimation-criteria}
给定参数函数 $g(\theta)$，以 $\widehat g=\widehat g(\mathbf X)$ 估计它；定义偏差、方差与\keyterm{均方误差}（mean squared error, MSE）。
在二阶矩存在时推导均方误差等于方差加偏差平方，说明无偏性并不自动意味着误差最小。
引入\keyterm{损失函数}（loss function）与\keyterm{风险函数}（risk function），明确比较范围及其对评价结果的影响。

\section{矩估计}
\label{sec:ps-method-moments}
介绍\keyterm{矩估计}（method of moments）：用样本矩替代总体矩并求解参数方程。
说明所用矩的存在、参数可识别性、解的存在唯一性与参数空间约束。
通过单参数和双参数模型比较计算简便性与统计效率，不把方程有解等同于解一定合理。

\section{最大似然估计}
\label{sec:ps-maximum-likelihood}
定义\keyterm{最大似然估计}（maximum likelihood estimation, MLE）为在 $\Theta$ 上最大化给定观测的似然。
讨论对数似然、边界解、多解、上确界不取到以及参数相关支撑的情形。
得分方程只用于寻找适当内点候选，最终需检查可行域与整体最大性；不变性在正确处理参数映射与解集后表述。

\section{无偏性与最小方差无偏估计}
\label{sec:ps-unbiased-estimation}
通过样本均值、样本方差与均匀分布端点估计比较无偏修正。
运用充分统计量、条件化与完备性构造\keyterm{一致最小方差无偏估计}（uniformly minimum-variance unbiased estimator, UMVUE）。
明确“最优”限于对所有参数均无偏的估计量类，不与所有可能损失下的最优性混同。

\section{相合性与渐近性质}
\label{sec:ps-estimation-asymptotics}
定义\keyterm{相合性}（consistency，又称一致性）与渐近正态性，区分它们与有限样本无偏性。
用大数定律与连续映射证明典型矩估计的相合性。
最大似然估计的渐近结论列出可识别性、内点真参数、适当光滑性和信息非退化等条件，并通过边界或不规则模型说明结论并非无条件成立。

\section{例题、习题与学习检查}
规划伯努利、泊松、正态和均匀端点模型的矩估计与最大似然估计；比较均匀端点估计的偏差和均方误差。
习题既包含计算，也包含“无偏但不相合”“有偏但相合”的例子与辨析。
学习检查要求分别报告估计目标、估计量公式、观测后的估计值和评价标准。


\chapter{区间估计与置信区域}
\label{chap:ps-intervals}

本章用区间或区域表达估计的不确定性，核心是重复抽样意义下的覆盖性质。
先修内容为抽样分布、点估计与中心极限定理；学习目标是区分精确区间与渐近区间，并正确解释置信水平。

\section{置信区间的定义与枢轴量}
\label{sec:ps-confidence-definition}
令 $0<\alpha<1$，定义\keyterm{置信区间}（confidence interval）$[L(\mathbf X),U(\mathbf X)]$ 的覆盖概率。
对精确的至少 $1-\alpha$ 水平，要求对每个 $\theta\in\Theta$ 有
\begin{equation}
  \mathbb P_\theta\{L(\mathbf X)\leq\theta\leq U(\mathbf X)\}\geq1-\alpha.
  \label{eq:ps-confidence-coverage}
\end{equation}
介绍\keyterm{枢轴量}（pivotal quantity）的构造与反解，区分随机区间和观测后的固定区间。
频率学派框架中的参数固定，不把置信水平直接解释为观测后参数落入该固定区间的概率。

\section{单个正态总体的参数区间}
\label{sec:ps-normal-intervals}
分别推导方差已知或未知时均值的双侧和单侧区间，以及正态总体方差的卡方区间。
明确正态性、独立同分布、样本量与自由度条件，统一分位数采用左侧累计概率的约定。
比较标准误、置信水平和区间长度之间的关系。

\section{两个总体的比较与配对数据}
\label{sec:ps-two-sample-intervals}
讨论独立正态样本的均值差区间，区分方差已知、未知但相等、未知且不等的情形。
给出合并方差的精确 $t$ 方法与韦尔奇近似（Welch approximation）的条件及区别。
配对样本先构造对内差值，再对独立的配对单位进行推断；方差比区间另需两个正态样本独立。

\section{比例、大样本区间与样本量规划}
\label{sec:ps-asymptotic-intervals}
用渐近正态性与标准误构造大样本区间，并以二项比例为例比较沃尔德区间（Wald interval）、威尔逊区间（Wilson interval）和精确反演区间。
讨论小样本、接近边界的参数与离散分布对覆盖率的影响；近似区间不承诺式\eqref{eq:ps-confidence-coverage}的有限样本覆盖。
根据目标误差和置信水平规划样本量，明确使用已知或预估方差的前提。

\section{例题、习题与学习检查}
规划正态均值与方差区间、两种测量方法的配对差区间、二项比例在小样本下的区间比较。
习题要求从枢轴量的不等式完整反解端点，检查单侧与双侧分位数。
置信区域作为向量参数的自然推广介绍；预测区间与均值区间的区别在第\ref{chap:ps-regression}章进一步展开。


\chapter{假设检验的原理与方法}
\label{chap:ps-testing}

本章研究如何在控制错误概率的前提下，根据数据评价关于总体的假设。
先修内容为抽样分布、似然和区间估计；学习目标是完整陈述检验问题、选择统计量、校准拒绝域并解释结论。

\section{假设、错误概率与功效}
\label{sec:ps-testing-framework}
定义\keyterm{原假设}（null hypothesis）$H_0:\theta\in\Theta_0$ 与\keyterm{备择假设}（alternative hypothesis）$H_1:\theta\in\Theta_1$，其中 $\Theta_0,\Theta_1$ 为不相交的参数子集。
区分简单假设、复合假设、单侧和双侧问题，定义检验统计量、拒绝域、第一类和第二类错误。
水平 $\alpha$ 的检验要求在原假设参数集合上拒绝概率的上确界不超过 $\alpha$；\keyterm{功效}（power）描述在备择参数处拒绝原假设的概率。

\section{显著性水平与 p 值}
\label{sec:ps-p-values}
解释\keyterm{p 值}（p-value）作为在指定检验及原假设校准方式下衡量观测结果极端程度的量。
在嵌套拒绝域框架下说明它与可拒绝的最小显著性水平之间的关系，离散和复合原假设情形需注明定义方式。
不把 p 值解释为原假设为真的概率，也不把未拒绝解释为证实原假设。

\section{正态总体与比例的常用检验}
\label{sec:ps-classical-tests}
推导单样本与双样本均值检验、配对检验、正态方差检验及独立正态样本方差比检验。
讨论二项比例的精确检验和大样本检验，明确每种方法的模型条件、统计量分布与自由度。
对同一问题比较检验结论与相应置信区间，建立匹配方法下的对偶关系。

\section{最强检验与似然比检验}
\label{sec:ps-likelihood-ratio-tests}
在简单原假设对简单备择假设的条件下陈述\keyterm{奈曼--皮尔逊引理}（Neyman--Pearson lemma），讨论离散情形达到指定水平可能需要随机化。
介绍\keyterm{似然比检验}（likelihood-ratio test），在约束参数空间与完整参数空间分别最大化似然。
渐近卡方结论仅在适当可识别、光滑、内点和信息非退化条件下使用，边界参数或不规则模型另行分析。

\section{检验设计、效应量与多重比较}
\label{sec:ps-power-multiplicity}
研究样本量、显著性水平、效应大小与功效之间的关系，区分统计显著与问题中的实际差异大小。
说明预先确定分析方案、报告区间和效应量的作用；介绍多次检验对总体错误概率的影响。
以邦费罗尼方法（Bonferroni method）说明一类基本控制方式，并明确所控制的是族错误率。

\section{例题、习题与学习检查}
规划均值是否超过给定阈值、两独立总体均值差、配对测量差、伯努利模型的最强检验。
习题要求写齐假设、条件、统计量、拒绝规则、观测结果与语境解释，并计算选定备择下的功效。
学习检查强调检验针对模型中的假设，数据质量与模型适用性不能由一个 p 值代替。


\chapter{拟合优度、列联表与非参数方法}
\label{chap:ps-nonparametric}

本章研究分布形状、类别关联及较少依赖具体分布形式的推断。
先修内容为统计检验与经验分布；学习目标是按数据类型和研究问题选择方法，并识别各方法实际检验的假设。

\section{卡方拟合优度检验}
\label{sec:ps-goodness-fit}
引入\keyterm{拟合优度检验}（goodness-of-fit test），由类别观测频数与假设期望频数构造皮尔逊卡方统计量（Pearson chi-square statistic）。
区分完全给定的类别概率与由同一数据估计参数的模型，说明自由度调整所需的正则条件。
讨论期望频数较小时近似可能失效、类别合并须有依据，以及精确或模拟校准的替代方法。

\section{列联表中的独立性与同质性}
\label{sec:ps-contingency}
介绍\keyterm{列联表}（contingency table），比较独立性检验与同质性检验的抽样机制和原假设。
推导期望频数、约束与自由度；对小样本二维列联表介绍费舍尔精确检验（Fisher's exact test）的条件化思想。
结合条件概率说明统计关联不能单独确立因果关系。

\section{基于经验分布的检验}
\label{sec:ps-empirical-tests}
介绍柯尔莫哥洛夫--斯米尔诺夫检验（Kolmogorov--Smirnov test）比较经验分布与理论分布的思路。
标准单样本分布无关校准以完全给定的连续原假设分布为前提；由样本估计参数或出现离散取值时须重新校准。
说明图形比较、检验统计量和显著性结论各自承担的任务。

\section{符号检验、秩检验与置换思想}
\label{sec:ps-rank-tests}
介绍\keyterm{非参数方法}（nonparametric methods）中的符号检验、威尔科克森符号秩检验（Wilcoxon signed-rank test）及秩和检验（rank-sum test）。
分别说明中位数、对称位置模型与两样本分布比较的假设，不能把秩和检验无条件解释为中位数相等检验。
引入\keyterm{置换检验}（permutation test），以原假设下的可交换性或随机分组机制说明置换的合法性，并处理并列值。

\section{例题、习题与学习检查}
规划多类别比例是否符合指定模型、类别变量是否独立、配对差的符号检验以及两独立样本的秩比较。
习题要求根据抽样设计说明检验对象和校准方法，并比较参数方法与非参数方法的条件。
学习检查强调“非参数”仍包含抽样、独立性、对称性或可交换性等假设。


\part{统计模型与数据分析}

\chapter{线性回归与预测}
\label{chap:ps-regression}

本章用线性模型研究响应变量如何随解释变量变化，并量化估计与预测误差。
先修内容为条件期望、协方差、估计和检验；多元部分需要矩阵秩、正交投影与二次型。

\section{回归问题与模型假设}
\label{sec:ps-regression-model}
定义\keyterm{响应变量}（response variable）、\keyterm{解释变量}（explanatory variable）与\keyterm{线性回归}（linear regression）。
先采用固定设计的一元模型 $Y_j=\beta_0+\beta_1x_j+\varepsilon_j$：$x_j$ 为给定设计值，$\beta_0,\beta_1$ 为未知系数，$\varepsilon_j$ 为随机误差。
分别说明零均值、同方差、不相关或独立、正态性在不同结论中的作用；随机设计情形改用给定解释变量的条件假设。

\section{最小二乘估计与几何解释}
\label{sec:ps-least-squares}
定义\keyterm{最小二乘法}（least squares），明确优化变量为回归系数、目标为观测残差平方和。
在设计值不全相同时推导一元回归解；多元模型以设计矩阵列满秩保证系数解唯一。
用正交投影解释拟合值与残差，区分线性代数优化事实与需要随机误差假设的统计结论。

\section{估计量的分布与高斯--马尔可夫定理}
\label{sec:ps-gauss-markov}
推导回归系数估计量的期望、协方差与误差方差估计。
在固定满秩设计、误差零均值及协方差为标量单位阵等条件下，陈述高斯--马尔可夫定理（Gauss--Markov theorem）的最佳线性无偏性。
精确 $t$ 与 $F$ 推断另加正态误差假设；最佳线性无偏不等于在所有估计量中最优。

\section{回归系数推断与预测区间}
\label{sec:ps-regression-inference}
构造斜率和一般线性系数的置信区间及检验，介绍整体回归显著性的 $F$ 检验。
区分给定解释变量处平均响应的区间与新观测的\keyterm{预测区间}（prediction interval），说明后者还包含新误差的波动。
预测新观测时明确其与训练误差的关系，并说明外推超出数据范围时依赖更强的模型可信度。

\section{诊断、多元回归与模型局限}
\label{sec:ps-regression-diagnostics}
介绍残差图、异常观测、异方差、共线性与影响点；多元模型说明“控制其他变量”的系数解释所依赖的模型。
定义决定系数并限定其解释，说明较高拟合程度不能单独证明预测性能或因果关系。
模型选择、稳健标准误与训练验证划分作为后续拓展，说明其解决的问题而不把软件输出当作自动结论。

\section{例题、习题与学习检查}
规划一元直线拟合、给定设计点的均值与预测区间、多元设计矩阵秩不足的例子。
推导题分别检查最小化、无偏性和区间覆盖所需条件；应用题要求把变量单位、估计系数与不确定性一起解释。
学习检查区分误差与残差、相关与回归、描述关系与因果解释。


\chapter{方差分析与试验设计初步}
\label{chap:ps-anova}

本章研究多个总体均值的比较，以及如何通过试验设计提高比较的可信度。
先修内容为正态抽样分布、假设检验与线性模型；学习目标是把模型、平方和分解和研究设计联系起来。

\section{单因素方差分析模型}
\label{sec:ps-oneway-anova}
建立\keyterm{单因素方差分析}（one-way analysis of variance, ANOVA）的固定效应模型。
区分组别、处理、观测单位与随机误差，明确独立性、共同误差方差及正态性用于精确检验的条件。
用基准组或效应和为零的约束处理参数表示的非唯一性。

\section{平方和分解与 F 检验}
\label{sec:ps-anova-decomposition}
将总变异分解为组间与组内变异，推导平方和、自由度和均方，构造各组均值相等的 $F$ 检验。
由线性模型正交分解解释正态误差下相关平方和的分布与独立性。
说明总体拒绝后只能断言并非所有均值都相等，不能直接确定具体哪一对不同。

\section{对比、多重比较与条件检查}
\label{sec:ps-anova-contrasts}
定义组均值的线性\keyterm{对比}（contrast），说明系数和为零的要求，区分事先规划与事后探索。
构造对比的估计与区间，结合第\ref{chap:ps-testing}章讨论多重比较控制。
通过残差检查讨论方差不齐、非正态和不平衡设计；替代方法需同时说明新的假设和推断目标。

\section{随机化、重复、区组与双因素设计}
\label{sec:ps-experimental-design}
介绍\keyterm{试验设计}（experimental design）的随机化、独立重复与\keyterm{区组}（block）原则。
区分技术重复与独立实验单位，说明伪重复会低估不确定性。
双因素设计中介绍主效应与交互作用；能否识别和检验交互作用取决于设计及重复情况，不能只看表格形式。

\section{例题、习题与学习检查}
规划三个处理组的均值比较、随机区组对照、存在交互作用的双因素示例。
习题要求先写模型再列方差分析表，核对各自由度之和，并说明实验单位及随机分组机制。
学习检查要求把显著性、差异大小、区间与设计限制一并解释。


\part{进一步学习的专题}

\chapter{贝叶斯统计初步（选读）}
\label{chap:ps-bayesian}

本章在统计模型中为未知参数引入概率分布，研究如何由先验与数据得到后验推断。
先修内容为贝叶斯公式、条件期望与似然；学习目标是理解推断所依赖的先验、模型和损失函数。

\section{先验、似然与后验}
\label{sec:ps-bayesian-model}
定义\keyterm{先验分布}（prior distribution）、\keyterm{后验分布}（posterior distribution）与边缘似然（marginal likelihood）。
在连续参数情形，用先验密度与似然的乘积构造后验，并检查归一化常数有限且非零。
说明这里关于参数的概率陈述属于所建立的贝叶斯模型，与频率学派固定参数的重复抽样解释不同。

\section{共轭模型与后验更新}
\label{sec:ps-conjugate-priors}
介绍\keyterm{共轭先验}（conjugate prior），推导贝塔--二项、伽马--泊松与方差已知的正态均值模型。
逐项注明超参数、数据统计量和更新后的参数，沿用前文伽马分布的率参数约定。
以顺序更新说明在模型条件一致时批量和逐批处理数据的联系。

\section{贝叶斯估计与可信区间}
\label{sec:ps-bayesian-decisions}
在后验期望损失最小的准则下介绍\keyterm{贝叶斯估计}（Bayes estimator）；平方损失和绝对损失分别联系后验均值与中位数。
介绍\keyterm{可信区间}（credible interval），并与第\ref{chap:ps-intervals}章的置信区间比较其概率解释。
最大后验估计单独说明为密度最大化方法，并讨论它依赖参数化的特点。

\section{后验预测与先验敏感性}
\label{sec:ps-posterior-predictive}
通过对参数的后验分布积分构造\keyterm{后验预测分布}（posterior predictive distribution），同时反映参数不确定性与新观测波动。
比较不同合理先验下的推断，说明不适当先验即使产生形式表达式也需检查后验是否为概率分布。
模型比较与贝叶斯因子仅作进一步学习入口，标明对先验和归一化常数的依赖。

\section{例题、习题与学习检查}
规划伯努利成功概率的顺序更新、泊松强度预测、正态均值估计的收缩现象。
习题要求分别报告先验、似然、后验、损失准则与预测对象。
学习检查避免把不同统计框架中的区间解释混用，并比较数据量与先验影响之间的关系。


\chapter{随机过程初步（选读）}
\label{chap:ps-processes}

本章把单个随机变量或有限随机向量推广为随时间变化的随机系统。
先修内容为联合分布、条件期望与极限定理；学习目标是掌握有限状态马尔可夫链和齐次泊松过程的基本模型。

\section{随机过程与时间结构}
\label{sec:ps-process-language}
定义\keyterm{随机过程}（stochastic process）、指标集、状态空间、有限维分布与样本路径。
区分离散时间和连续时间、离散状态和连续状态，并通过随机游走说明时间相关性。
只给出均值和协方差通常不能确定整个过程的分布。

\section{有限状态马尔可夫链}
\label{sec:ps-markov-chains}
定义\keyterm{马尔可夫性质}（Markov property）与时间齐次转移矩阵，说明初始分布也是模型的一部分。
推导多步转移概率与查普曼--柯尔莫哥洛夫方程（Chapman--Kolmogorov equations）。
介绍可达、互通、闭类、吸收状态、常返与暂留以及周期的含义，主线限定在有限状态空间。

\section{平稳分布、极限分布与长期平均}
\label{sec:ps-stationarity}
定义\keyterm{平稳分布}（stationary distribution）并求解归一化线性方程。
说明有限不可约链具有唯一平稳分布；从任意初始分布收敛到它还需非周期性。
用周期链区分平稳分布与边缘分布极限，并说明有限不可约链的长期时间平均结论不要求非周期性。
介绍细致平衡与可逆性，为第\ref{chap:ps-computation}章的模拟方法作准备。

\section{泊松过程与等待时间}
\label{sec:ps-poisson-process}
定义强度为 $\lambda>0$ 的\keyterm{齐次泊松过程}（homogeneous Poisson process），说明起点、独立增量、平稳增量及计数分布。
研究计数与到达时刻的联系，推导独立指数间隔和伽马到达时间分布。
讨论独立泊松过程叠加及独立标记下的稀疏化，明确这些独立性条件不可省略。

\section{例题、习题与学习检查}
规划天气状态的简化链、带吸收边界的随机游走、泊松到达的等待与计数问题。
习题要求写明初始分布和转移规则、检查平稳方程，并构造平稳分布存在但分布不收敛的周期例子。
鞅与布朗运动仅列为后续课程方向，不作为本书主干的先修要求。


\chapter{随机模拟与统计计算（选读）}
\label{chap:ps-computation}

本章用算法近似概率量和统计推断结果，建立理论对象、计算过程与误差评价之间的联系。
先修内容为极限定理与统计推断；马尔可夫链方法另需第\ref{chap:ps-processes}章。
每种算法的正文均明确输入、输出、初始化、更新、停止规则及适用条件。

\section{随机数与分布抽样}
\label{sec:ps-random-generation}
区分数学上的独立随机样本与计算机\keyterm{伪随机数}（pseudorandom numbers），记录种子和生成方式以便复现实验。
讲解逆分布函数抽样与拒绝抽样；后者要求可计算的目标与提议密度关系，以及有效的包络常数。
说明实际算法输出的是有限精度数值，理论分布性质需要结合实现方式理解。

\section{蒙特卡罗积分与误差估计}
\label{sec:ps-monte-carlo}
介绍\keyterm{蒙特卡罗方法}（Monte Carlo method），将积分写成期望，再以独立样本均值估计。
有限绝对一阶矩用于大数定律保证；有限非零方差用于通常的中心极限定理误差近似。
以预设模拟次数或有理论依据的序贯规则停止，不能不加说明地反复检查普通固定样本区间后自行停止。

\section{重要性抽样与方差缩减}
\label{sec:ps-variance-reduction}
介绍\keyterm{重要性抽样}（importance sampling）、控制变量与对偶变量的基本思想。
重要性抽样要求提议分布覆盖目标积分的有关支撑，并检查加权变量的可积性与二阶矩。
区分普通加权估计与自归一化估计的有限样本性质；比较方法时同时报告精度与计算成本。

\section{自助法与置换计算}
\label{sec:ps-bootstrap}
介绍\keyterm{自助法}（bootstrap），从经验分布重复抽样近似统计量的抽样分布，计算标准误与区间。
普通非参数自助法以适当的独立同分布抽样和统计量正则性为基础；极值、边界或依赖数据需要另行论证。
区分自助重抽样与原假设下的置换，两者的目标和合法性依据不同；有限次重复还引入模拟误差。

\section{马尔可夫链蒙特卡罗}
\label{sec:ps-mcmc}
介绍\keyterm{马尔可夫链蒙特卡罗}（Markov chain Monte Carlo, MCMC），先在有限状态情形构造以目标分布为平稳分布的链。
给出梅特罗波利斯--黑斯廷斯算法（Metropolis--Hastings algorithm）与吉布斯抽样（Gibbs sampling）的更新思想。
分别讨论不变分布、不可约性、非周期性与时间平均收敛所需条件；连续状态推广需补充相应理论。
相关输出不能直接当作独立样本，诊断图与有效样本量只提供计算证据，不能单独证明已经收敛。

\section{例题、习题与学习检查}
规划用模拟核对概率计算、比较两种积分估计方差、模拟检验功效与区间覆盖率，以及用简单离散目标分布检查马尔可夫链抽样。
每个实验同时报告模型、重复次数、统计量、随机性来源和模拟标准误；与可解析的小模型作交叉核对。
学习检查区分模型误差、抽样误差、蒙特卡罗误差与数值优化误差。


\appendix

\chapter{必要的数学工具}
\label{app:ps-prerequisites}

本附录按正文需求补充工具，每个主题以可直接使用的定义、条件、公式和短例题组织，不承担完整先修课程的替代作用。

\section{集合、计数与常用求和}
整理集合运算、可数集合、排列组合、二项式定理、几何级数与常见求和恒等式。
对无限级数区分绝对收敛与条件收敛，说明重排及逐项运算需要条件。

\section{微积分与积分交换}
回顾广义积分、多元积分、变量替换、泰勒展开及余项控制。
对非负或绝对可积函数说明交换求和与积分的适用情形；单调收敛、控制收敛与富比尼定理的概率应用列为进阶工具，并准确陈述使用条件。

\section{线性代数与二次型}
回顾矩阵秩、逆、迹、行列式、正交投影、正定与半正定二次型。
解释协方差矩阵与最小二乘问题使用这些工具的位置，补充实对称矩阵的正交对角化。

\section{特殊积分与渐近记号}
整理伽马函数、贝塔函数及其与常见密度归一化的关系，说明参数范围。
引入确定性及概率意义的阶记号时分别定义，配合极限定理说明近似项与余项的区别。


\chapter{查阅表、方法索引与综合学习检查}
\label{app:ps-reference}

本附录的表格与索引将在各章正文完成后同步整理，所有条目回指推导所在章节，并保留假设，避免形成脱离条件的公式清单。

\section{常见分布查阅表}
按分布名称、参数及范围、支撑、质量函数或密度、分布函数、均值、方差和典型生成机制组织。
统一几何及负二项的计数约定、伽马的率参数约定、正态的方差参数约定；存在多种参数化时列出转换关系。

\section{估计、区间与检验方法索引}
按推断目标、抽样机制、模型条件、所用统计量、精确或渐近分布以及正文位置组织。
将独立样本、配对样本、有限总体抽样分开列示，并标注边界参数和小样本情形需要重新检查的方法。

\section{综合例题与分层习题安排}
规划三条贯穿全书的综合问题：伯努利成功概率从分布到推断，正态测量误差从均值方差到回归，计数与等待从泊松分布到泊松过程。
各章习题依次设置概念辨析、基本计算、证明反例和综合建模；选定习题提供完整解答，其余提供提示或结果并明确标识。

\section{完成一项概率统计分析的检查顺序}
先确认研究问题、随机对象与数据获得方式，再写出模型及假设、确定目标量和方法、核对适用条件、执行推导或计算，最后解释结果与不确定性。
正文中的每个综合案例均按这一顺序组织，使分散的概念能够汇成可独立完成的分析过程。


\end{document}
